\SourceLeaf{122}
\Article{150}
Formae hae plures habent proprietates satis memorabiles, quarum tamen unam tantummodo apponimus. Si $B$ est numerus compositus ad $A$ primus, inter cuius factores primos occurrunt $2m$, qui in aliqua forma non-divisorum ipsius $xx-A$ continentur, $B$ in aliqua forma divisorum ipsius $xx-A$ contentus erit; si vero multitudo factorum primorum ipsius $B$ in aliqua forma non-divisorum ipsius $xx-A$ contentorum impar est, $B$ quoque in forma non-divisorum contentus erit. Demonstrationem quae non est difficilis omittimus. Hinc vero sequitur, non modo quemvis numerum primum sed etiam quemvis compositum imparem ad $A$ primum, qui in aliqua forma non-divisorum contineatur, non-divisorem fore: necessario enim aliquis factor primus talis numeri debet esse non-divisor.

\begin{center}
\textit{De aliorum laboribus circa has investigationes.}
\end{center}

\Article{151}
Theorema fundamentale. quod sane inter elegantissima in hoc genere est referendum, in eadem forma simplici, in qua supra propositum est, a nemine hucusque fuit prolatum Quod eo magis est mirandum. quum aliae quaedam propositiones illi superstruendae, ex quibus ad illud facile reveniri potuisset, ill. Eulero iam innotuerint. Formas certas dari, in quibus omnes divisores primi numerorum formae $xx-A$ contineantur, aliasque in quibus omnes non-divisores primi numerorum eiusdem formae sint comprehensi, ita ut hae illas excludant, noverat methodumque illas formas inveniendi eruerat: sed omnes ipsius conatus ad demonstrationem perveniendi semper irriti fuerunt, veritatique illi per inductionem inventae maiorem tantummodo verisimilitudinem conciliaverunt. In aliqua quidem tractatione, \textit{Novae demonstrationes circa divisores numerorum formae $xx+nyy$}, quae in Acad. Petrop. recitata est 1775 Nov. 20, et post mortem viri summi in \textit{T.} I. \textit{Nov. Act.} huius Ac. p. 47 sqq. est conservata, voti se compotem credidisse videtur: sed hic error irrepsit, scilicet p. 65 \textit{tacite} supposuit, formas tales divisorum et non-divisorum exstare\HistoricalNote{*}{Nempe dari numeros $r,r',r''$ etc.; $n,n',n''$ etc. omnes diversos et $<4A$ tales ut omnes divisores primi ipsius $xx-A$ sub aliqua formarum $4Ak+r$, $4Ak+r'$ etc. contineantur. omnesque non-divisores primi sub aliqua harum $4Ak+n$, $4Ak+n'$ etc (designante $k$ numerum indeterminatum).}, unde non difficile erat \textit{quales} esse debeant derivare: methodus autem qua usus est ad comprobationem illius suppositionis haud
\SourceLeaf{123}
idonea videtur. In alio schediasmate, \textit{De criteriis aequationis $fxx+gyy=hzz$ utrumque resolutionem admittat necne}, Opusc. Anal. T. I. (ubi $f,g,h$ sunt dati, $x,y,z$ indeterminati) per inductionem invenit, si aequatio pro aliquo valore ipsius $h=s$ solubilis sit, eandem pro quovis alio valore ipsi $s$ secundum mod. $4fg$ congruo, siquidem sit numerus primus, solubilem fore, ex qua propositione suppositio de qua diximus haud difficile demonstrari potest. Sed etiam huius theorematis demonstratio omnes ipsius labores elusit\HistoricalNote{*)}{Uti ipse fatetur, l. c. p. 216: ,,Huius elegantissimi theorematis demonstratio adhuc desideratur, postquam a pluribus iamdudum frustra est investigata.\,\ldots\ Quocirca plurimum is praestitisse censendus erit, cui successerit demonstrationem huius theorematis invenire.`` --- Quanto ardore vir immortalis demonstrationem huius theorematis aliorumque, quae tantummodo casus speciales theor. fundam. sunt. desideraverit, videre licet ex multis aliis locis Opusc. Anall. Conf. \textit{Additamentum ad diss.} VIII, T. I. et \textit{diss.} XIII, T. II. pluresque diss. in Comment. Petrop., iam passim laudatae.}, quod non est mirandum, quia nostro iudicio a theoremate fundamentali erat proficiscendum. Ceterum veritas huius propositionis ex iis quae in Sect. sequenti docebimus sponte demanabit.

Post Eulerum, clar. Le Gendre eidem argumento operam navavit, in egregia tract. \textit{Recherches d'analyse indéterminée}, Hist. de l'Ac. des Sc. 1785 p. 465 sqq., ubi pervenit ad theorema, quod si rem ipsam spectas cum th. fund. idem est, scilicet designantibus $p,q$ duos numeros primos positivos, fore residua absolute minima potestatum $p^{\frac{q-1}{2}}$, $q^{\frac{p-1}{2}}$ sec. mod. $q,p$ resp. aut ambo $+1$, aut ambo $-1$. quando aut $p$ aut $q$ sit formae $4n+1$; quando vero tum $p$ tum $q$ sit formae $4n+3$, alterum res. min. fore $+1$, alterum $-1$, p. 516. ex quo sec. art. 106 derivatur, \textit{relationem} (in signif. art. 146 acceptam) ipsius $p$ ad $q$ ipsiusque $q$ ad $p$ \textit{eandem} esse, quando aut $p$ aut $q$ sit formae $4n+1$, \textit{oppositam}, quando tum $p$ tum $q$ sit formae $4n+3$. Propos. haec inter propp. art. 131 est contenta, sequitur etiam ex 1, 3, 9, art. 133; vicissim autem theor. fund. ex ipsa derivari potest. Clar. Le Gendre etiam demonstrationem tentavit, de qua quum perquam ingeniosa sit in Sect. seq. fusius loquemur. Sed quoniam in ea plura sine demonstratione supposuit (uti ipse fatetur p. 520 \textit{Nous avons supposé seulement} etc.), quae partim a nemine hucusque sunt demonstrata, partim nostro quidem iudicio sine theor. fund. ipso demonstrari nequeunt: via quam ingressus est, ad scopum deducere non posse videtur, nostraque demonstratio pro prima erit habenda. --- Ceterum infra \textit{duas alias demonstrationes} eiusdem gravissimi theorematis trademus, a praec. et inter se toto coelo diversas.

\SourceLeaf{124}
\begin{center}
\textit{De congruentiis secundi gradus non puris.}
\end{center}

\Article{152}
Hactenus congruentiam puram $xx\equiv A\ (\mathrm{mod.}\,m)$ tractavimus, ipsiusque resolubilitatem dignoscere docuimus. \textit{Radicum ipsarum} investigatio per art. 105 ad eum casum est reducta, ubi $m$ est aut primus aut primi potestas, posterior vero per art. 101 ad eum, ubi $m$ est primus. Pro hoc autem casu ea quae in art. 61 sqq. tradidimus una cum iis quae in Sectt. V et VIII docebimus, omnia fere complectuntur quae per methodos directas erui possunt. Sed hae ubi sunt applicabiles plerumque infinities prolixiores sunt quam indirectae quas in Sect. VI docebimus, adeoque non tam propter utilitatem suam in praxi quam propter pulcritudinem memorabiles. --- \textit{Congruentiae secundi gradus non purae} ad puras facile reduci possunt. Proposita congruentia
\[
axx+bx+c\equiv 0
\]
secundum mod. $m$ solvenda, huic aequivalebit congruentia
\[
4aaxx+4abx+4ac\equiv 0\ (\mathrm{mod.}\,4am)
\]
\textit{i. e.} quivis numerus alteri satisfaciens etiam alteri satisfaciet. Haec vero ita exhiberi potest
\[
(2ax+b)^2\equiv bb-4ac\ (\mathrm{mod.}\,4am)
\]
unde omnes valores ipsius $2ax+b$ minores quam $4am$ si qui dantur inveniri possunt. Quibus per $r,r',r''$ etc. designatis, omnes solutiones congr. prop. deducentur ex solutionibus congruentiarum
\[
2ax\equiv r-b,\quad 2ax\equiv r'-b\text{ etc.}\ (\mathrm{mod.}\,4am)
\]
quas in Sect. II invenire docuimus. Ceterum observamus, solutionem plerumque per varia artificia contrahi posse, \textit{ex. gr.} loco congr. prop. aliam inveniri posse
\[
a'xx+2b'x+c'\equiv 0
\]
illi aequipollentem, et in qua $a'$ ipsum $m$ metiatur; haec vero de quibus Sect. ultima conferri potest, hic explicare brevitas non permittit.

\BookRule{.35\linewidth}
\clearpage
\pdfbookmark[0]{Sectio quinta}{section-five}
\SourceLeaf{125}
\begin{center}
SECTIO QUINTA

DE

FORMIS AEQUATIONIBUSQUE INDETERMINATIS

SECUNDI GRADUS.
\end{center}
\BookRule{.35\linewidth}

\begin{center}
\textit{Disquisitionis propositum; formarum definitio et signum.}
\end{center}

\Article{153}
In hac sectione imprimis de functionibus duarum indeterminatarum $x,y$, huius formae
\[
axx+2bxy+cyy
\]
ubi $a,b,c$ sunt integri dati, tractabimus, quas \textit{formas secundi gradus} sive simpliciter \textit{formas} dicemus. Huic disquisitioni superstruetur solutio problematis famosi. invenire omnes solutiones aequationis cuiuscunque indeterminatae secundi gradus duas incognitas implicantis, sive hae incognitae valores integros sive rationales tantum nancisci debeant. Problema hoc quidem iam ab ill. La Grange in omni generalitate est solutum, multaque insuper ad naturam \textit{formarum} pertinentia tum ab hoc ipso magno geometra tum ab ill. Eulero partim primum inventa, partim, a Fermatio olim inventa, demonstrationibus munita. Sed nobis acriter formarum perquisitioni insistentibus tam multa nova se obtulerunt, ut totum argumentum ab integro resumere operae pretium duxerimus, eo magis, quod Virorum illorum inventa, multis locis sparsa, paucis innotuisse experti sumus; porro quod methodus per quam haec tractabimus nobis ad maximam partem est propria; tandem
\SourceLeaf{126}
quod nostra sine nova illorum expositione ne intelligi quidem possent. Nullum vero dubium nobis esse videtur, quin multa eaque egregia in hoc genere adhuc lateant. in quibus alii vires suas exercere possint. Ceterum quae ad veritatum insignium historiam pertinent, loco suo semper trademus.

Formam $axx+2bxy+cyy$, quando de indeterminatis $x,y$ non agitur, ita designabimus, $(a,b,c)$. Haec itaque expressio denotabit indefinite summam trium partium, producti numeri dati $a$ in quadratum indeterminatae cuiuscunque; producti duplicati numeri $b$ in hanc indeterminatam in aliam indeterminatam; producti numeri $c$ in quadratum huius secundae indeterminatae. \textit{Ex. gr.} $(1,0,2)$ exprimet summam quadrati et quadrati duplicati. Ceterum, quamvis formae $(a,b,c)$ et $(c,b,a)$ idem designent, si ad \textit{partes ipsas} tantum respicimus, tamen different si insuper ad partium \textit{ordinem} attendimus: quare sedulo eas in posterum distinguemus; quid vero inde lucremur in sequentibus sufficienter patebit.

\Needspace{8\baselineskip}
\begin{center}
\textit{Numerorum repraesentatio; determinans.}
\end{center}

\Article{154}
Numerum aliquem datum per formam datam \textit{repraesentari} dicemus, si formae indeterminatis tales valores integri tribuuntur, ut ipsius valor numero dato fiat aequalis. Hic habebimus sequens

\textsc{Theorema.} \textit{Si numerus $M$ ita per formam $(a,b,c)$ repraesentari potest, ut indeterminatarum valores, per quos hoc efficitur, inter se sint primi: erit $bb-ac$ residuum quadraticum numeri $M$.}

\textit{Dem.} Sint valores indeterminatarum $m,n$, scilicet
\[
amm+2bmn+cnn=M
\]
accipianturque numeri $\mu,\nu$ ita ut sit $\mu m+\nu n=1$ (art. 40). Tum per evolutionem facile probatur esse
\[
\begin{aligned}
&amm+2bmn+cnn)(a\nu\nu-2b\mu\nu+c\mu\mu\\
&\qquad={}(\mu(mb+nc)-\nu(ma+nb))^2-(bb-ac)(m\mu+n\nu)^2
\end{aligned}
\]
sive
\[
M(a\nu\nu-2b\mu\nu+c\mu\mu)=(\mu(mb+nc)-\nu(ma+nb))^2-(bb-ac)
\]
Quare erit
\SourceLeaf{127}
\[
bb-ac\equiv(\mu(mb+nc)-\nu(ma+nb))^2\ (\mathrm{mod.}\,M)
\]
\textit{i. e.} $bb-ac$ residuum quadraticum ipsius $M$.

Numerum $bb-ac$, a cuius indole proprietates formae $(a,b,c)$ imprimis pendere in sequentibus docebimus, \textit{determinantem} huius formae vocabimus.

\begin{center}
\textit{Valores expr. $\surd(bb-ac)(\mathit{mod.}\,M)$ ad quos repraesentatio numeri $M$ par formam $(a,b,c)$ pertinet.}
\end{center}

\Article{155}
Erit itaque
\[
\mu(mb+nc)-\nu(ma+nb)
\]
valor expressionis
\[
\surd(bb-ac)\ (\mathrm{mod.}\,M)
\]
Constat autem, numeros $\mu,\nu$ infinitis modis ita determinari posse ut sit $\mu m+\nu n=1$, unde alii aliique valores illius expressionis prodibunt, qui quem nexum inter se habeant videamus. Sit non modo $\mu m+\nu n=1$, sed etiam $\mu'm+\nu'n=1$, ponaturque
\[
\mu(mb+nc)-\nu(ma+nb)=v,\quad
\mu'(mb+nc)-\nu'(ma+nb)=v'
\]
Multiplicando aequationem $\mu m+\nu n=1$ per $\mu'$, alteram $\mu'm+\nu'n=1$ per $\mu$, et subtrahendo fit $\mu'-\mu=n(\mu'\nu-\mu\nu')$ similiterque multiplicando illam per $\nu'$, hanc per $\nu$, fit subtrahendo $\nu'-\nu=m(\mu\nu'-\mu'\nu)$. Hinc statim prodit
\[
v'-v=(\mu'\nu-\mu\nu')(amm+2bmn+cnn)=(\mu'\nu-\mu\nu')M
\]
sive $v'\equiv v\ (\mathrm{mod.}\,M)$. Quomodocunque igitur $\mu,\nu$ determinentur, formula $\mu(mb+nc)-\nu(ma+nb)$ valores \textit{diversos} (\textit{i. e.} incongruos) expressionis $\surd(bb-ac)\ (\mathrm{mod.}\,M)$ dare nequit. Si itaque $v$ est valor quicunque illius formulae: repraesentationem numeri $M$ per formam $axx+2bxy+cyy$ eam ubi $x=m$, $y=n$, \textit{pertinere} dicemus \textit{ad valorem $v$ expressionis $\surd(bb-ac)(\mathit{mod.}\,M)$}. Ceterum facile ostendi potest, si valor formulae illius aliquis sit $v$ atque $v'\equiv v\ (\mathrm{mod.}\,M)$, loco numerorum $\mu,\nu$, qui dant $v$, alios $\mu',\nu'$ accipi posse, qui dent $v'$. Scilicet faciendo
\[
\mu'=\mu+\frac{n(v'-v)}{M},\quad
\nu'=\nu-\frac{m(v'-v)}{M}
\]
fiet
\[
\mu'm+\nu'n=\mu m+\nu n=1
\]
\SourceLeaf{128}
valor autem formulae ex $\mu',\nu'$ prodiens superabit valorem ex $\mu,\nu$ prodeuntem quantitate $(\mu'\nu-\mu\nu')M$, quae fit $=(\mu m+\nu n)(v'-v)=v'-v$, sive valor ille erit $=v'$.

\Article{156}
Si duae repraesentationes eiusdem numeri $M$ per eandem formam $(a,b,c)$ habentur, in quibus indeterminatae valores inter se primos habent: hae vel ad eundem valorem expr. $\surd(bb-ac)\ (\mathrm{mod.}\,M)$ pertinere possunt vel ad diversos. Sit
\[
M=amm+2bmn+cnn=am'm'+2bm'n'+cn'n'
\]
atque
\[
\mu m+\nu n=1,\quad \mu'm'+\nu'n'=1
\]
patetque si fuerit
\[
\mu(mb+nc)-\nu(ma+nb)\equiv\mu'(m'b+n'c)-\nu'(m'a+n'b)\ (\mathrm{mod.}\,M)
\]
congruentiam semper manere, quicunque alii valores idonei pro $\mu,\nu$; $\mu',\nu'$ accipiantur, in quo casu utramque repraesentationem ad \textit{eundem} valorem expr. $\surd(bb-ac)\ (\mathrm{mod.}\,M)$ pertinere dicemus; si vero congruentia pro ullis valoribus ipsorum $\mu,\nu$; $\mu',\nu'$ locum non habet, pro nullis locum habebit, repraesentationesque ad valores \textit{diversos} pertinebunt. Si vero
\[
\mu(mb+nc)-\nu(ma+nb)\equiv-(\mu'(m'b+n'c)-\nu'(m'a+n'b))
\]
repraesentationes ad valores \textit{oppositos} expr. $\surd(bb-ac)$ pertinere dicentur. Omnibus hisce denominationibus etiam utemur, quando de pluribus repraesentationibus eiusdem numeri per formas \textit{diversas}, sed quae eundem determinantem habent, agitur.

\textit{Ex.} Sit forma proposita haec $(3,7,-8)$ cuius determinans $=73$. Per hanc formam habentur repraesentationes numeri $57$ hae
\[
3\mathbin{.}13^2+14\mathbin{.}13\mathbin{.}25-8\mathbin{.}25^2;\quad
3\mathbin{.}5^2+14\mathbin{.}5\mathbin{.}9-8\mathbin{.}9^2
\]
Pro prima poni potest $\mu=2$, $\nu=-1$, unde prodit valor expr. $\surd73\ (\mathrm{mod.}\,57)$ ad quam repr. pertinet
\[
{}=2(13\mathbin{.}7-25\mathbin{.}8)+(13\mathbin{.}3+25\mathbin{.}7)=-4
\]
\SourceLeaf{129}
Simili modo repraesentatio secunda pertinere invenitur, faciendo $\mu=2$, $\nu=-1$, ad valorem $+4$. Quare ambae repraesentationes ad valores oppositos pertinent.

Antequam ulterius progredimur, observamus, formas quarum determinans $=0$ ab investigationibus sequentibus prorsus exclusas esse, quippe quae theorematum concinnitatem tantummodo turbarent, adeoque tractationem peculiarem postulent.

\begin{center}
\textit{Forma aliam implicans, sive sub alia contenta; transformatio, propria et impropria.}
\end{center}

\Article{157}
Si forma $F$, cuius indeterminatae sunt $x,y$, in aliam $F'$, cuius indeterminatae sunt $x',y'$, per substitutiones tales
\[
x=\alpha x'+\beta y',\quad y=\gamma x'+\delta y'
\]
transmutari potest, ita ut $\alpha,\beta,\gamma,\delta$ sint integri: priorem \textit{implicare} posteriorem, sive posteriorem \textit{sub priori contentam esse} dicemus. Sit forma $F$ haec
\[
axx+2bxy+cyy
\]
forma $F'$ vero haec
\[
a'x'x'+2b'x'y'+c'y'y'
\]
habebunturque sequentes tres aequationes:
\[
\begin{aligned}
a'&=a\alpha\alpha+2b\alpha\gamma+c\gamma\gamma\\
b'&=a\alpha\beta+b(\alpha\delta+\beta\gamma)+c\gamma\delta\\
c'&=a\beta\beta+2b\beta\delta+c\delta\delta
\end{aligned}
\]
Multiplicando aequationem secundam per se ipsam, primam per tertiam, et subtrahendo fit deletis partibus se destruentibus
\[
b'b'-a'c'=(bb-ac)(\alpha\delta-\beta\gamma)^2
\]
Unde sequitur determinantem formae $F'$ per determinantem formae $F$ divisibilem et quotientem esse quadratum; manifesto igitur hi determinantes \textit{eadem signa} habebunt. Quodsi itaque insuper forma $F'$ per similem substitutionem in formam $F$ transmutari potest, \textit{i. e.} si tum $F'$ sub $F$, tum $F$ sub $F'$ contenta est,
\SourceLeaf{130}
formarum determinantes erunt aequales\HistoricalNote{*)}{Manifestum est ex analysi praecedente hanc propositionem etiam ad formas quarum determinans $=0$, patere. Sed aequatio $(\alpha\delta-\beta\gamma)^2=1$ ad hunc casum non est extendenda.} atque $(\alpha\delta-\beta\gamma)^2=1$. In hoc casu formas \textit{aequivalentes} dicemus. Quare ad formarum aequivalentiam aequalitas determinantium est conditio necessaria. licet illa ex hac sola minime sequatur. --- Substitutionem $x=\alpha x'+\beta y'$, $y=\gamma x'+\delta y'$ vocabimus \textit{transformationem propriam}, si $\alpha\delta-\beta\gamma$ est numerus positivus, \textit{impropriam}, si $\alpha\delta-\beta\gamma$ est negativus; formam $F'$ \textit{proprie} aut \textit{improprie} sub forma $F$ contentam esse dicemus, si $F$ per transformationem propriam aut impropriam in formam $F'$ transmutari potest. Si itaque formae $F,F'$ sunt aequivalentes, erit $(\alpha\delta-\beta\gamma)^2=1$, adeoque si transformatio est propria, $\alpha\delta-\beta\gamma=+1$, si est impropria, $=-1$. --- Si plures transformationes simul sunt propriae, aut simul impropriae, \textit{similes} eas dicemus; propriam contra et impropriam \textit{dissimiles}.

\begin{center}
\textit{Aequivalentia, propria et impropria.}
\end{center}

\Article{158}
\textit{Si formarum $F,F'$ determinantes sunt aequales atque $F'$ sub $F$ contenta: etiam $F$ sub $F'$ contenta erit et quidem proprie vel improprie prout $F'$ sub $F$ proprie vel improprie continetur.} Transeat $F$ in $F'$ ponendo
\[
x=\alpha x'+\beta y',\quad y=\gamma x'+\delta y'
\]
transibitque $F'$ in $F$ ponendo
\[
x'=\delta x-\beta y,\quad y'=-\gamma x+\alpha y
\]
Patet enim per hanc substitutionem ex $F'$ fieri idem, quod fiat ex $F$ ponendo
\[
x=\alpha(\delta x-\beta y)+\beta(-\gamma x+\alpha y),\quad
y=\gamma(\delta x-\beta y)+\delta(-\gamma x+\alpha y)
\]
sive
\[
x=(\alpha\delta-\beta\gamma)x,\quad y=(\alpha\delta-\beta\gamma)y
\]
Hinc vero manifesto ex $F$ fit $(\alpha\delta-\beta\gamma)^2F$ \textit{i. e.} rursus $F$ (art. praec.). Perspicuum autem est, transformationem posteriorem esse propriam vel impropriam, prout prior sit propria vel impropria.

Si tum $F'$ sub $F$, tum $F$ sub $F'$ \textit{proprie} continentur, formas \textit{proprie aequi\SourceLeaf{131}valentes}, si illae sub invicem improprie, vocabimus \textit{improprie aequivalentes}. --- Ceterum usus harum distinctionum mox innotescet.

\textit{Exempl.} Forma $2xx-8xy+3yy$ per substitutiones $x=2x'+y'$, $y=3x'+2y'$ transit in formam $-13x'x'-12x'y'-2y'y'$, haec vero in illam factis $x'=2x-y$, $y'=-3x+2y$. Quare formae $(2,-4,3)$, $(-13,-6,-2)$ erunt \textit{proprie aequivalentes}.

Problemata quae tractare iam aggrediemur sunt haec: I. Propositis duabus formis quibuscunque eundem determinantem habentibus. investigare utrum sint aequivalentes necne, utrum proprie aut improprie aut utroque modo, nam etiam hoc fieri potest. Quando vero determinantes inaequales habent, annon saltem altera alteram implicet, proprie vel improprie vel utroque modo. Denique invenire omnes transformationes alterius in alteram, tam proprias quam improprias. II. Proposita forma quacunque, invenire utrum numerus datus per eam repraesentari possit omnesque repraesentationes assignare. Sed quoniam formae determinantis negativi hic aliam methodum requirunt quam formae determinantis positivi, primo trademus ea quae utrisque sunt communia, tum vero formas cuiusvis generis seorsim considerabimus.

\begin{center}
\textit{Formae oppositae.}
\end{center}

\Article{159}
\textit{Si forma $F$ formam $F'$ implicat, haec vero formam $F''$, forma $F$ etiam formam $F''$ implicabit.}

Sint indeterminatae formarum $F,F',F''$ respective $x,y$; $x',y'$; $x'',y''$ transeatque $F$ in $F'$ ponendo
\[
x=\alpha x'+\beta y'.\quad y=\gamma x'+\delta y'
\]
$F'$ in $F''$ ponendo
\[
x'=\alpha'x''+\beta'y'',\quad y'=\gamma'x''+\delta'y''
\]
patetque, $F$ in $F''$ transmutatum iri ponendo
\[
x=\alpha(\alpha'x''+\beta'y'')+\beta(\gamma'x''+\delta'y''),\quad
y=\gamma(\alpha'x''+\beta'y'')+\delta(\gamma'x''+\delta'y'')
\]
sive
\[
x=(\alpha\alpha'+\beta\gamma')x''+(\alpha\beta'+\beta\delta')y'',\quad
y=(\gamma\alpha'+\delta\gamma')x''+(\gamma\beta'+\delta\delta')y''
\]
Quare $F$ ipsam $F''$ implicabit.
